\documentclass[10pt,twocolumn]{article}

\usepackage[margin=0.75in,columnsep=0.28in]{geometry}
\usepackage{amsmath,amssymb,mathtools}
\usepackage{booktabs}
\usepackage{enumitem}
\usepackage{microtype}
\usepackage{xcolor}
\usepackage[colorlinks=true,linkcolor=blue!55!black,citecolor=blue!55!black,urlcolor=blue!55!black]{hyperref}

\newcommand{\R}{\mathbb{R}}
\newcommand{\diag}{\operatorname{diag}}
\newcommand{\blkdiag}{\operatorname{blkdiag}}
\newcommand{\half}{\tfrac{1}{2}}
\newcommand{\ones}{\mathbf{1}}
\DeclareMathOperator*{\argmin}{arg\,min}

\setlength{\abovedisplayskip}{4pt plus 1pt minus 1pt}
\setlength{\belowdisplayskip}{4pt plus 1pt minus 1pt}
\emergencystretch=1.5em

\title{Forward and Backward Passes of a Log-Barrier ADMM QP Solver\\
\large A perturbed-KKT derivation in the TurboMPC ADMM convention}
\author{Jianghan Zhang}
\date{\today}

\begin{document}
\maketitle

\section{From trajectory optimization to a QP}
\label{sec:trajopt}

We consider the discrete-time optimal control problem (OCP) in the form
of~\cite{turbompc},
\begin{equation}
\label{eq:ocp}
\begin{aligned}
\min_{(x_t,u_t)_{t=0}^{N}}\ & \textstyle\sum_{t=0}^{N-1}\ell_t(x_t,u_t)+\ell_N(x_N)\\
\text{s.t.}\ \ & x_0 = x_{\mathrm{init}},\\
& f_t(x_t,u_t,x_{t+1},u_{t+1}) = 0, \quad t=0,\dots,N{-}1,\\
& \underline g_t \le g_t(x_t,u_t) \le \overline g_t, \quad\ \ t=0,\dots,N,
\end{aligned}
\end{equation}
with states $x_t\in\R^{n_x}$, controls $u_t\in\R^{n_u}$, implicit discretized dynamics
$f_t$, and pointwise-in-time inequality constraints $g_t$ (bounds, obstacles,
friction cones, \dots). The OCP data depend on parameters $\theta$ (cost weights,
constraint geometry) with respect to which we ultimately want gradients.

\paragraph{Stacked NLP.} Overloading $x := (x_t,u_t)_{t=0}^{N}\in\R^{n_w}$ for the
stacked stage variables, collect the cost into $\ell$, the initial-condition and
dynamics residuals into $f\in\R^{m}$, and the inequalities into \emph{one-sided}
rows $g\in\R^{p}$ (each finite bound of \eqref{eq:ocp} contributes one row,
$g_t \le \overline g_t$ or $-g_t \le -\underline g_t$; infinite bounds are dropped),
so that \eqref{eq:ocp} is the parametric nonlinear program
\begin{equation}
\label{eq:nlp}
\min_{x}\ \ell(x;\theta)\quad\text{s.t.}\quad f(x;\theta) = 0,\quad
g(x;\theta) \le h(\theta).
\end{equation}

\paragraph{SQP.} Sequential quadratic programming~\cite[Ch.~18]{nw} solves
\eqref{eq:nlp} by iterating: at the current primal--dual iterate, quadratize the
Lagrangian, linearize the constraints, solve the resulting QP, and update the
iterate along its solution with a linesearch. With the constraint Jacobians
$C := \partial f/\partial x$ and $G := \partial g/\partial x$ evaluated at the
iterate, $P \approx \nabla^2_{xx}\mathcal L$ a Hessian approximation (in practice
Gauss--Newton), and $q$ the corresponding gradient, each subproblem is the convex QP
studied in the rest of this note:
\begin{equation}
\label{eq:qp}
\mathrm{QP}(\theta):\quad
\min_{x}\ \half x^\top P x + q^\top x
\quad\text{s.t.}\quad Cx = c,\ \ Gx \le h,
\end{equation}
with $P\succeq 0$, $C\in\R^{m\times n_w}$, $G\in\R^{p\times n_w}$. The time
structure of \eqref{eq:ocp} makes $P$ block-tridiagonal, $C$ block-bidiagonal, and
$G$ block-diagonal~\cite{turbompc} --- structure both passes below exploit. A
differentiable MPC layer must provide (i) a \emph{forward pass} producing the
primal--dual solution of \eqref{eq:qp} at each SQP iteration
(Section~\ref{sec:forward}) and (ii) a \emph{backward pass} producing
$\nabla_\theta L$ for a scalar loss $L$ of the converged solution --- obtained by
differentiating the \emph{NLP's} optimality system, not the QP's
(Section~\ref{sec:backward}).

\section{Perturbed KKT and smoothed complementarity}
\label{sec:pkkt}

With multipliers $y_f\in\R^m$ (equalities) and $y_g\in\R^p$ (inequalities) and the
slack $s := h - Gx$, the KKT conditions of \eqref{eq:qp} are
\begin{equation}
\label{eq:kkt}
\begin{aligned}
&Px+q+C^\top y_f+G^\top y_g = 0, \qquad Cx = c,\\
&Gx + s = h,\qquad s\ge 0,\quad y_g\ge 0,\quad s\circ y_g = 0,
\end{aligned}
\end{equation}
where $\circ$ is the elementwise product. All difficulty --- combinatorial for the
solver, and geometric for differentiation --- sits in the last line. The solution map
$\theta \mapsto (x,y_f,y_g)$ of \eqref{eq:kkt} is only piecewise smooth: at points
where \emph{strict complementarity} fails ($s_i = y_{g,i} = 0$ for some row) the KKT
Jacobian is singular and the map is nondifferentiable, with gradient jumps across
active-set changes~\cite{optnet,frey}. We therefore smooth \eqref{eq:kkt} with two
perturbations, used together or separately.

\paragraph{(i) Barrier perturbation $\kappa$.} Following the interior-point
tradition, replace complementarity by the \emph{perturbed KKT} (central-path)
condition~\cite[Ch.~19]{nw}
\begin{equation}
\label{eq:centralpath}
s \circ y_g = \kappa\ones, \qquad s > 0,\ y_g > 0,
\end{equation}
for a barrier parameter $\kappa > 0$. This is the stationarity system of the
equality-constrained log-barrier subproblem (add $-\kappa\sum_i\log s_i$ to the
cost); as $\kappa\searrow 0$ it recovers \eqref{eq:kkt}. Differentiating the
$\kappa$-perturbed system instead of \eqref{eq:kkt} is exactly the mechanism used to
obtain smooth gradients through QP/NLP layers~\cite{tracy,frey}.

\paragraph{(ii) Inequality-slack (Moreau) perturbation $(\delta_\xi,\gamma_\xi)$.}
Following TurboMPC's slack convention~\cite{turbompc}, introduce an
\emph{inequality slack} $\xi \ge 0$ per one-sided row, toggled by the binary
indicator $\delta_\xi\in\{0,1\}$ and priced quadratically:
$Gx - \delta_\xi\xi \le h$ with penalty $\delta_\xi\tfrac{\gamma_\xi}{2}\|\xi\|^2$,
$\gamma_\xi > 0$. For $\delta_\xi = 1$, minimizing over $\xi$ first replaces the
indicator of $\{Gx\le h\}$ by its Moreau envelope
$\tfrac{\gamma_\xi}{2}\,\mathrm{dist}^2(Gx,\,\{\cdot\le h\})$~\cite{parikh}: the
constraint wall becomes a stiff quadratic ramp, the QP stays feasible under SQP
linearization error, and the multiplier ramps up continuously as
$y_g = \gamma_\xi\xi$ (TurboMPC's \texttt{ADMMSlack} smoothing~\cite{turbompc}).
\textbf{The default is $\delta_\xi = 0$} (no inequality slack);
$(\delta_\xi{=}1,\,\gamma_\xi{\to}\infty)$ is equivalent to $\delta_\xi = 0$.

\paragraph{Combined perturbed system.} Applying both, the log-barrier subproblem is
\begin{equation}
\label{eq:barrier-sub}
\begin{aligned}
\min_{x,\xi}\ & \half x^\top Px + q^\top x
+ \delta_\xi\,\frac{\gamma_\xi}{2}\|\xi\|^2
 - \kappa\sum_{i=1}^{p}\log s_i\\
\text{s.t.}\ & Cx = c, \qquad s = h - Gx + \delta_\xi\,\xi ,
\end{aligned}
\end{equation}
whose stationarity in $s$ gives $y_g = \kappa/s$ (elementwise) and, when
$\delta_\xi = 1$, stationarity in $\xi$ gives $\gamma_\xi\xi = y_g$. Eliminating
$\xi = \gamma_\xi^{-1}y_g$, the optimality conditions collapse to a square smooth
root-finding problem in $w := (x, y_f, y_g)$:
\begin{equation}
\label{eq:pkkt}
F_{\mathrm{QP}}(w;\theta) :=
\begin{bmatrix}
Px + q + C^\top y_f + G^\top y_g\\[1pt]
Cx - c\\[1pt]
s \circ y_g - \kappa\ones
\end{bmatrix} = 0,
\end{equation}
\begin{equation}
\label{eq:sdef}
s := h - Gx + \delta_\xi\gamma_\xi^{-1}y_g \;(>0), \qquad y_g > 0 .
\end{equation}
The forward pass converges to the root of \eqref{eq:pkkt} for each QP subproblem;
the backward pass differentiates the nonlinear counterpart of this system
(Section~\ref{sec:backward}). The two perturbations place the solver in a four-way
family of ADMM methods:
\begin{center}
\small
\begin{tabular}{@{}lll@{}}
\toprule
& $\kappa = 0$ (hard compl.) & $\kappa > 0$ (central path)\\
\midrule
$\delta_\xi=0$ & KKT \eqref{eq:kkt}; OSQP~\cite{osqp} & pure barrier (\textbf{default})\\
$\delta_\xi=1$ & Moreau; \texttt{ADMMSlack}~\cite{turbompc} & slack barrier\\
\bottomrule
\end{tabular}
\end{center}
Both $\kappa > 0$ cells are modes of this solver (Section~\ref{sec:forward}).
For any $\kappa>0$ the residual $F_{\mathrm{QP}}$ is smooth and its Jacobian is
nonsingular (Section~\ref{sec:backward}), so $\theta\mapsto w(\theta)$ is $C^1$ ---
including across activation/deactivation of constraints. The price is an
$O(\kappa)$ bias of $x(\theta)$ relative to the hard QP solution, and an
$O(\|y_g\|/\gamma_\xi)$ constraint violation from the inequality slack when
$\delta_\xi = 1$; $\kappa\to 0$ (with $\delta_\xi = 0$, or $\gamma_\xi\to\infty$)
recovers \eqref{eq:kkt}.

\section{Forward pass: ADMM with a log-barrier slack update}
\label{sec:forward}

We follow the ADMM convention of \cite{turbompc} (itself OSQP's
splitting~\cite{osqp,boyd}). Stack the constraints as
$A := \bigl[\begin{smallmatrix}C\\ G\end{smallmatrix}\bigr]$ and give $Ax$ its own
copy $z = (z_f, z_g)$:
\begin{equation}
\label{eq:split}
\begin{aligned}
\min_{x,z,\xi}\ & \half x^\top Px + q^\top x
+ \delta_\xi\,\frac{\gamma_\xi}{2}\|\xi\|^2
+ \mathcal I_{\{c\}}(z_f) + B_\kappa(z_g,\xi)\\
\text{s.t.}\ & Ax - z = 0,
\end{aligned}
\end{equation}
where $\mathcal I$ denotes an indicator function and
\begin{equation}
\label{eq:Bkappa}
B_\kappa(z_g,\xi) :=
-\kappa\sum_{i=1}^{p}\log\bigl(h_i - z_{g,i} + \delta_\xi\,\xi_i\bigr)
\end{equation}
is the log barrier on the copy. (At $\kappa = 0$, $B_\kappa$ degenerates to the
indicator of $\{z_g - \delta_\xi\xi \le h\}$ and \eqref{eq:split} is exactly the
TurboMPC splitting.) Introducing duplicated variables
$(\tilde x, \tilde z)$ as in \cite{turbompc}, multipliers $w$ for
$\tilde x - x = 0$ (which vanishes identically after the first iteration and is
dropped below) and $y = (y_f, y_g)$ for $\tilde z - z = 0$, and step-size parameters
$\sigma > 0$, $\rho := \blkdiag(\rho_f I_m,\ \rho_g I_p)\succ 0$, the augmented
Lagrangian is
\begin{multline}
\label{eq:auglag}
\mathcal L_{\sigma,\rho}
= \half\tilde x^\top P\tilde x + q^\top\tilde x
+ \delta_\xi\,\frac{\gamma_\xi}{2}\|\xi\|^2
+ \mathcal I_{\{c\}}(z_f) + B_\kappa(z_g,\xi)\\
+ \mathcal I_{\{0\}}(A\tilde x - \tilde z)
+ \frac{\sigma}{2}\bigl\|\tilde x - x\bigr\|^2
+ \half\bigl\|\tilde z - z + \rho^{-1}y\bigr\|_\rho^2 ,
\end{multline}
with $\|v\|_\rho^2 := v^\top\rho\, v$. One ADMM iteration minimizes
$\mathcal L_{\sigma,\rho}$ over $(\tilde x,\tilde z)$, then over $(x, z, \xi)$, then
ascends the duals, with over-relaxation $\alpha\in(0,2)$.

\subsection{Primal update}
Minimizing \eqref{eq:auglag} over $(\tilde x,\tilde z)$ is an equality-constrained
strictly convex QP whose KKT conditions are the linear
system~\cite[Eq.~(7)]{turbompc}
\begin{equation}
\label{eq:primal}
\begin{bmatrix} P+\sigma I & A^\top\\ A & -\rho^{-1}\end{bmatrix}
\begin{bmatrix}\tilde x^{k+1}\\ \nu^{k+1}\end{bmatrix}
=
\begin{bmatrix}\sigma x^{k} - q\\ z^{k} - \rho^{-1}y^{k}\end{bmatrix}.
\end{equation}
The $-\rho^{-1}$ block is eliminated by a Schur complement:
\begin{equation}
\label{eq:schur}
S\,\tilde x^{k+1} = \sigma x^{k} - q + A^\top\!\bigl(\rho z^{k} - y^{k}\bigr),
\qquad \tilde z^{k+1} = A\tilde x^{k+1},
\end{equation}
with $S := P + \sigma I + A^\top\!\rho A \succ 0$. By the structure noted in
Section~\ref{sec:trajopt}, $S$ is block-tridiagonal in time; its sparse
factorization is computed once and reused across iterations, refactorizing only when
$\rho$ is adapted~\cite{turbompc,osqp}. The over-relaxed iterates are
\begin{equation}
\label{eq:overrelax}
\hat x^{k+1} = \alpha\tilde x^{k+1} + (1-\alpha)x^{k},\quad
\hat z^{k+1} = \alpha\tilde z^{k+1} + (1-\alpha)z^{k},
\end{equation}
and the $x$-minimization gives $x^{k+1} = \hat x^{k+1}$.

\subsection{Slack update: the proximal operator of the log barrier}
\label{sec:slackupdate}
Minimizing \eqref{eq:auglag} over $z_f$ gives $z_f^{k+1} = c$ (so $z_f$ need not be
stored). The inequality block is the essential step. Let
\begin{equation}
\label{eq:ztilde}
\tilde z_g^{k+1} := \hat z_g^{k+1} + \rho_g^{-1}y_g^{k}
\end{equation}
be the usual ADMM target --- the same point OSQP projects onto $\{z_g\le h\}$ and
TurboMPC's \texttt{ADMMSlack} projects smoothly. Here the update is the prox of the
barrier, jointly over $(z_g,\xi)$ and separable across rows:
\begin{equation}
\label{eq:zprox}
(z_g,\xi)^{k+1}
= \argmin_{z,\xi}\, B_\kappa(z,\xi)
+ \delta_\xi\tfrac{\gamma_\xi}{2}\|\xi\|^2
+ \tfrac{\rho_g}{2}\|z - \tilde z_g^{k+1}\|^2
\end{equation}
(for $\delta_\xi = 0$ the inequality slack is fixed, $\xi \equiv 0$, exactly as in
TurboMPC's slack update). Per row, with the barrier slack
$s(z,\xi) := h - z + \delta_\xi\xi$, the objective of \eqref{eq:zprox} is
\begin{equation}
\label{eq:prox-obj}
\varphi(z,\xi) = -\kappa\log s(z,\xi)
+ \delta_\xi\tfrac{\gamma_\xi}{2}\xi^2
+ \tfrac{\rho_g}{2}\bigl(z - \tilde z_g^{k+1}\bigr)^2 .
\end{equation}
This is an \emph{unconstrained} smooth problem: the barrier absorbs $s \ge 0$ into
its domain ($\varphi\to+\infty$ as $s\to 0^+$), and the quadratic terms make
$\varphi$ strictly convex and coercive there, so the unique minimizer is the
zero-gradient point. With $\partial s/\partial z = -1$ and
$\partial s/\partial\xi = \delta_\xi$, the chain rule gives
\begin{equation}
\label{eq:prox-stat}
\begin{aligned}
\partial_z\varphi
&= \tfrac{\kappa}{s} + \rho_g\bigl(z - \tilde z_g^{k+1}\bigr) = 0
&&\Rightarrow\
\tfrac{\kappa}{s} = \rho_g\bigl(\tilde z_g^{k+1} - z\bigr),\\
\partial_\xi\varphi
&= -\tfrac{\kappa}{s} + \gamma_\xi\xi = 0\ \ (\delta_\xi{=}1)
&&\Rightarrow\
\gamma_\xi\xi = \tfrac{\kappa}{s},
\end{aligned}
\end{equation}
and the stationary $\xi = \kappa/(\gamma_\xi s) > 0$ is automatically positive, so
$\xi \ge 0$ never activates. Solving \eqref{eq:prox-stat} for
$z = \tilde z_g^{k+1} - \kappa/(\rho_g s)$ and
$\delta_\xi\xi = \delta_\xi\kappa/(\gamma_\xi s)$ and substituting back into
$s = h - z + \delta_\xi\xi$ with the residual $r^{k+1} := h - \tilde z_g^{k+1}$
yields $s = r^{k+1} + \Gamma/s$, i.e.\ the scalar quadratic
$s^2 - r^{k+1}s - \Gamma = 0$ per row, whose positive root is, elementwise,
\begin{equation}
\label{eq:retr-update}
\boxed{
\begin{aligned}
s^{k+1} &= b_\Gamma\bigl(r^{k+1}\bigr),
& \Gamma &:= \kappa\bigl(\rho_g^{-1} + \delta_\xi\,\gamma_\xi^{-1}\bigr),\\
\xi^{k+1} &= \delta_\xi\,\frac{\kappa}{\gamma_\xi\, s^{k+1}},
& z_g^{k+1} &= h - s^{k+1} + \xi^{k+1},
\end{aligned}}
\end{equation}
where $b_\mu$ is the \emph{retraction map} (positive root of $b^2 - vb - \mu = 0$)
\begin{equation}
\label{eq:b}
b_\mu(v) := \frac{v + \sqrt{v^2 + 4\mu}}{2}, \qquad \mu > 0,
\end{equation}
with the elementary properties:
\begin{enumerate}[nosep,leftmargin=2.2em,label=(\roman*)]
\item $b_\mu(v) > 0$ for all $v\in\R$;
\item $b_\mu(v)\,b_\mu(-v) = \mu$ (built-in complementarity);
\item $b_\mu(v) - b_\mu(-v) = v$;
\item $b_\mu(v)\to\max(v,0)$ as $\mu\to 0^+$ (smooth ReLU);
\item $b_\mu'(v) = \half\bigl(1+v/\sqrt{v^2+4\mu}\bigr)\in(0,1)$ and
  $b_\mu'(v) + b_\mu'(-v) = 1$.
\end{enumerate}
In the default mode $\delta_\xi = 0$: $\Gamma = \kappa/\rho_g$, $\xi^{k+1} = 0$, and
$z_g^{k+1} = h - b_{\kappa/\rho_g}(r^{k+1})$. (For $v<0$ we evaluate
$b_\mu(v) = 2\mu/(\sqrt{v^2+4\mu}-v)$ to avoid numerical issues.)

\subsection{Dual update and exact perturbed complementarity}
The dual ascent uses the over-relaxed copy~\cite{turbompc}:
\begin{equation}
\label{eq:dual}
y_f^{k+1} = y_f^{k} + \rho_f\bigl(\hat z_f^{k+1} - c\bigr),\quad
y_g^{k+1} = y_g^{k} + \rho_g\bigl(\hat z_g^{k+1} - z_g^{k+1}\bigr).
\end{equation}
By \eqref{eq:ztilde}, $y_g^{k+1} = \rho_g(\tilde z_g^{k+1} - z_g^{k+1})$, and the
$z$-stationarity in \eqref{eq:prox-stat} evaluates this in closed form:
\begin{equation}
\label{eq:headline}
\boxed{\,
y_g^{k+1} = \frac{\kappa}{s^{k+1}}
\;\Rightarrow\;
s^{k+1}\!\circ y_g^{k+1} = \kappa\ones,\
\gamma_\xi\xi^{k+1} = \delta_\xi y_g^{k+1}
\,}
\end{equation}
elementwise, \emph{at every iteration $k$} --- not only in the limit. The iterates
live on the central-path manifold \eqref{eq:centralpath} by construction, with
$y_g^{k+1} > 0$ automatic from property (i), and the inequality-slack stationarity
holds identically in both modes. At a fixed point, $Ax = z$ makes
$s = h - Gx + \delta_\xi\xi = h - Gx + \delta_\xi\gamma_\xi^{-1}y_g$, so
$(x, y_f, y_g)$ satisfies \emph{exactly} the perturbed KKT system
\eqref{eq:pkkt}--\eqref{eq:sdef} of the QP.

\paragraph{Limits.} As $\kappa\to 0^+$, property (iv) collapses
\eqref{eq:retr-update}, for $\delta_\xi = 1$, to TurboMPC's smoothed projection
$\tilde\Pi^{\gamma_\xi}(\tilde z_g) =
\frac{\rho_g}{\gamma_\xi+\rho_g}\tilde z_g
+ \frac{\gamma_\xi}{\gamma_\xi+\rho_g}\Pi(\tilde z_g)$
\cite[Eq.~(15)]{turbompc}, and, for $\delta_\xi = 0$, to OSQP's hard projection
$z_g^{k+1} = \min(\tilde z_g^{k+1}, h)$: a smooth deformation of the standard
projection ADMM, at identical per-iteration cost.

\paragraph{Termination, $\rho$-adaptation, continuation.} Convergence is declared on
the standard residuals $r_{\mathrm{prim}} = \|Ax - z\|_\infty$ and
$r_{\mathrm{dual}} = \|Px + q + A^\top y\|_\infty$
\cite[Eq.~(3)]{turbompc} (slack stationarity
$\gamma_\xi\xi - \delta_\xi y_g = 0$ holds identically by \eqref{eq:headline});
$\bar\rho$ is adapted OSQP-style from the residual ratio
with $\rho_f = 10^3\bar\rho$, $\rho_g = \bar\rho$, $\sigma = 10^{-6}$, refactorizing
$S$ on each change. Optionally $\kappa$ is annealed geometrically with warm starts;
for differentiable use a fixed $\kappa$ is preferred so the backward differentiates
the fixed point the forward reached.

\section{Backward pass: implicit differentiation of the NLP's perturbed KKT}
\label{sec:backward}

Let $L = L(x^\star)$ be a scalar loss of the converged SQP solution. We do
\emph{not} differentiate through solver iterations --- neither the ADMM loop nor the
SQP loop. Nor do we differentiate the last QP subproblem's optimality system: the
object that carries the sensitivities of \eqref{eq:nlp} is the \emph{NLP's} own
smoothed KKT system, and QP-level sensitivities coincide with the NLP's only when
the QP carries the \emph{exact} Lagrangian Hessian at the
solution~\cite[Thm.~2]{frey}. Since the forward SQP typically runs a Gauss--Newton
$P$ (Section~\ref{sec:trajopt}), differentiating the QP it actually solved can
silently degrade the gradients~\cite[Rem.~3]{frey}. The backward pass therefore
re-assembles the NLP system at the converged point.

\subsection{The NLP's smoothed KKT system}
At SQP convergence, the iterate $w^\star := (x^\star, y_f, y_g)$ satisfies the
$(\kappa,\delta_\xi,\gamma_\xi)$-perturbed optimality conditions of \eqref{eq:nlp}
--- the nonlinear counterpart of \eqref{eq:pkkt}, with the smoothed complementarity
imposed on the \emph{nonlinear} residual (the interior-point-smoothed conditions of
\cite[Eq.~(10)]{frey}, whose barrier parameter $\tau$ is our $\kappa$):
\begin{equation}
\label{eq:nlp-pkkt}
F(w;\theta) :=
\begin{bmatrix}
\nabla_x\ell(x;\theta) + C(x)^\top y_f + G(x)^\top y_g\\[1pt]
f(x;\theta)\\[1pt]
s \circ y_g - \kappa\ones
\end{bmatrix} = 0,
\end{equation}
\begin{equation}
\label{eq:nlp-sdef}
s := h(\theta) - g(x;\theta) + \delta_\xi\gamma_\xi^{-1}y_g \;(>0),
\qquad y_g > 0,
\end{equation}
where $C(x) := \partial f/\partial x$ and $G(x) := \partial g/\partial x$ are the
constraint Jacobians evaluated at $x$ (for affine $f,g$ and quadratic $\ell$,
\eqref{eq:nlp-pkkt} reduces exactly to \eqref{eq:pkkt}). The first row is
$\nabla_x$ of the NLP Lagrangian
$\mathcal L := \ell(x;\theta) + y_f^\top f(x;\theta) + y_g^\top g(x;\theta)$.
Differentiating \eqref{eq:nlp-pkkt} at its root yields the \emph{exact} gradient of
the relaxed solution map, smooth across active-set changes.

\subsection{The implicit function theorem}
The residual $F(w;\theta)$ of \eqref{eq:nlp-pkkt} is $C^1$ in $(w,\theta)$ whenever
$\ell, f, g$ are twice continuously differentiable~\cite[Ass.~1]{frey}. Let
$w^\star$ satisfy $F(w^\star;\theta) = 0$ with $s > 0$, $y_g > 0$ (guaranteed for
$\kappa>0$ by \eqref{eq:headline}; at SQP convergence the linearized and nonlinear
residuals coincide), and define
$M := \partial F/\partial w\,\big|_{(w^\star,\theta)}$. If $M$ is nonsingular
(established below), the implicit function theorem yields a locally unique $C^1$
map $\theta\mapsto w(\theta)$ with $F(w(\theta);\theta)\equiv 0$ and, by total
differentiation
$\;\frac{\partial F}{\partial w}\frac{\partial w}{\partial\theta}
+ \frac{\partial F}{\partial\theta} = 0$,
\begin{equation}
\label{eq:ift}
\frac{\partial w}{\partial\theta} = -\,M^{-1}\frac{\partial F}{\partial\theta}
\end{equation}
--- the sensitivity result of \cite[Thm.~1]{frey}; their Thm.~3 adds that the
smoothed map is $C^1$ for small $\kappa$ and approaches the unsmoothed one with
error $O(\kappa)$. Rather than forming $\partial w/\partial\theta$
column-by-column ($n_\theta$ solves with $M$), the gradient of a scalar loss needs
a \emph{single} solve with $M^\top$: absorbing the sign, define the adjoint
$\lambda = (\lambda_x,\lambda_f,\lambda_g)$ by
\begin{equation}
\label{eq:adjoint}
M^\top \lambda = -\bigl(\nabla_x L,\, 0,\, 0\bigr),
\qquad
\nabla_\theta L = \Bigl(\frac{\partial F}{\partial\theta}\Bigr)^{\!\top}\lambda .
\end{equation}

\subsection{The Jacobian $M$ and the adjoint system}
Differentiating \eqref{eq:nlp-pkkt} with $\partial s/\partial x = -G(x)$ and
$\partial s/\partial y_g = \delta_\xi\gamma_\xi^{-1}I$ gives, in block form (rows =
residuals, columns = $x, y_f, y_g$; all quantities evaluated at $w^\star$),
\begin{equation}
\label{eq:M}
M =
\begin{bmatrix}
H & C^\top & G^\top\\
C & 0 & 0\\
-\diag(y_g)\,G & 0 & \diag\!\bigl(s + \delta_\xi\gamma_\xi^{-1}y_g\bigr)
\end{bmatrix},
\end{equation}
where the $(1,1)$ block is the \emph{exact Lagrangian Hessian} of the NLP,
\begin{equation}
\label{eq:H}
H := \nabla^2_{xx}\mathcal L
= \nabla^2_{xx}\ell
+ \sum_{j=1}^{m} y_{f,j}\,\nabla^2_{xx} f_j
+ \sum_{i=1}^{p} y_{g,i}\,\nabla^2_{xx} g_i ,
\end{equation}
\emph{not} the Gauss--Newton $P$ used by the forward SQP --- sensitivities computed
with an approximate Hessian can be badly degraded even when the forward solution is
accurate~\cite[Thm.~2, Rem.~3]{frey}. Because $f$ and $g$ are stagewise, $H$ retains
$P$'s block-tridiagonal sparsity. The adjoint system \eqref{eq:adjoint} written out
is
\begin{subequations}
\label{eq:adjeqs}
\begin{align}
H\lambda_x + C^\top\lambda_f - G^\top\diag(y_g)\,\lambda_g &= -\nabla_x L,
\label{eq:adj-a}\\
C\lambda_x &= 0, \label{eq:adj-b}\\
G\lambda_x + \diag\!\bigl(s+\delta_\xi\gamma_\xi^{-1}y_g\bigr)\lambda_g &= 0.
\label{eq:adj-c}
\end{align}
\end{subequations}
Two equivalent reductions of \eqref{eq:adjeqs} are used in practice.

\subsection{Reduction A: retraction parameterization}
\label{sec:redA}
On the central-path manifold (stated for the default $\delta_\xi = 0$), properties
(ii)--(iii) of $b_\kappa$ make the pair $(s, y_g)$ the graph of a \emph{single}
unconstrained variable $v := y_g - s\in\R^p$:
\begin{equation}
s = b_\kappa(-v), \qquad y_g = b_\kappa(v),
\end{equation}
which satisfies $s\circ y_g = \kappa\ones$, $s,y_g>0$ \emph{identically} --- the
complementarity residual disappears from the system. Differentiating with
$B_p := \diag\bigl(b_\kappa'(v)\bigr)$, $B_n := \diag\bigl(b_\kappa'(-v)\bigr)$
(both positive by (v)) gives $dy_g = B_p\,dv$, $ds = -B_n\,dv$, and the remaining
residuals (stationarity, primal inequality $g(x;\theta) + s - h$, equality)
linearize in $(dx, dv, dy_f)$ with matrix
$\bigl[\begin{smallmatrix} H & G^\top B_p & C^\top\\ G & -B_n & 0\\
C & 0 & 0\end{smallmatrix}\bigr]$.
Subtracting $G^\top\times$(second row) from the first and using
$B_p + B_n = I$ (property (v)) symmetrizes it:
\begin{equation}
\label{eq:J}
J =
\begin{bmatrix}
H - G^\top G & G^\top & C^\top\\
G & -B_n & 0\\
C & 0 & 0
\end{bmatrix} = J^\top .
\end{equation}
Because $J$ is symmetric, the adjoint solve is
\begin{equation}
\label{eq:J-adjoint}
J\begin{bmatrix}\lambda_x\\ \lambda_v\\ \lambda_f\end{bmatrix}
= -\begin{bmatrix}\nabla_x L\\ 0\\ 0\end{bmatrix},
\qquad
\bar\lambda_g := B_p\,\lambda_v ,
\end{equation}
where $\bar\lambda_g$ is the effective inequality adjoint used in the gradient
assembly below.

\subsection{Reduction B: Schur elimination onto a weighted Hessian}
\label{sec:redB}
Alternatively --- and for either value of $\delta_\xi$ --- eliminate $\lambda_g$
directly from \eqref{eq:adj-c}: its coefficient
$\diag(s+\delta_\xi\gamma_\xi^{-1}y_g)$ is diagonal and strictly positive for
$\kappa>0$, so
\begin{equation}
\label{eq:lamg}
\lambda_g = -\diag\!\bigl(s+\delta_\xi\gamma_\xi^{-1}y_g\bigr)^{-1} G\,\lambda_x .
\end{equation}
Substituting into \eqref{eq:adj-a} folds the inequalities into the Hessian block:
\begin{equation}
\label{eq:Wsys}
\begin{bmatrix}
H + G^\top\!\diag(W)\,G & C^\top\\
C & 0
\end{bmatrix}
\begin{bmatrix}\lambda_x\\ \lambda_f\end{bmatrix}
= -\begin{bmatrix}\nabla_x L\\ 0\end{bmatrix},
\end{equation}
with the smooth per-row weight
\begin{equation}
\label{eq:W}
W := \frac{y_g}{s + \delta_\xi\gamma_\xi^{-1}y_g},
\qquad
W\big|_{\delta_\xi=0} = \frac{y_g}{s} = \frac{y_g^2}{\kappa},
\end{equation}
and the effective inequality adjoint
$\bar\lambda_g := -y_g\circ\lambda_g = \diag(W)\,G\lambda_x$, matching
\eqref{eq:J-adjoint}. This reduced saddle system is exactly the symmetric reduced
coefficient matrix of \cite[Eq.~(9)]{frey} (there $S^{-1}M = \diag(\mu/s)$ plays
the role of $W$). $W$ is the $C^1$ analogue of a hard active-set mask:
\begin{center}
\footnotesize
\begin{tabular}{@{}llll@{}}
\toprule
regime & $y_{g,i}$ & $s_i$ & $W_i$\\
\midrule
inactive row & $\to 0$ & $>0$ & $\to 0$ (row drops out)\\
active, $\kappa{\to}0$, $\delta_\xi{=}1$ & $>0$ & $\to\gamma_\xi^{-1}y_{g,i}$ & $\to\gamma_\xi$ (penalty stiffness)\\
active, $\kappa{\to}0$, $\delta_\xi{=}0$ & $>0$ & $\to 0$ & $\to\infty$ (hard constraint)\\
\bottomrule
\end{tabular}
\end{center}
so $W\in[0,\gamma_\xi]$ always when $\delta_\xi = 1$. Since $G$ is block-diagonal
per stage, $G^\top\diag(W)G$ perturbs only the diagonal blocks of $H$:
\eqref{eq:Wsys} has \emph{the same block-tridiagonal saddle structure as the forward
Schur system} \eqref{eq:schur}, and the same structured sparse solver serves both
passes~\cite{turbompc}.

\paragraph{Equivalence of the reductions.} Eliminating
$\lambda_v = B_n^{-1}G\lambda_x$ from the middle row of \eqref{eq:J} gives
$H - G^\top G + G^\top B_n^{-1}G = H + G^\top B_pB_n^{-1}G$ (using $I - B_n = B_p$),
and on the manifold
$B_pB_n^{-1} = \diag\bigl(b_\kappa'(v)/b_\kappa'(-v)\bigr)
= \diag\bigl(b_\kappa(v)/b_\kappa(-v)\bigr) = \diag(y_g/s)$,
which is exactly $W|_{\delta_\xi=0}$; likewise
$\bar\lambda_g = B_pB_n^{-1}G\lambda_x = \diag(W)G\lambda_x$. (For $\delta_\xi = 1$
the middle block becomes $-(B_n + \gamma_\xi^{-1}B_p)$ and
$B_p(B_n+\gamma_\xi^{-1}B_p)^{-1}
= \diag\bigl(y_g/(s+\gamma_\xi^{-1}y_g)\bigr) = W$.)
Reductions A and B are the same linear system.

\paragraph{Nonsingularity --- why $\kappa$ is load-bearing.}
\eqref{eq:Wsys} is nonsingular iff $C$ has full row rank (LICQ; true by
construction: the initial-condition and dynamics rows) and
$H + G^\top\diag(W)G \succ 0$ on $\ker C$, which for small $\kappa$ follows from
the second-order sufficient conditions at the solution~\cite[Ass.~1, Thm.~1]{frey}.
The elimination \eqref{eq:lamg} itself required
$\diag(s + \delta_\xi\gamma_\xi^{-1}y_g)\succ 0$: for $\kappa>0$ this holds
\emph{everywhere}, because $s\circ y_g = \kappa$ keeps $s$ away from zero. In the
hard limit $\kappa=0$ the genuine failure is a \emph{weakly active} row,
$s_i = 0$ \emph{and} $y_{g,i} = 0$ (strict complementarity fails): the $i$-th
complementarity row of $M$ vanishes identically, $M$ is singular, and the solution
map is nondifferentiable~\cite{optnet,frey}, with $W_i = 0/0$ jumping between $0$
(inactive side) and $\gamma_\xi$ or $\infty$ (active side). (A strongly active row,
$s_i = 0$, $y_{g,i} > 0$, also zeroes the block entry when $\delta_\xi = 0$, but
harmlessly: $W_i\to\infty$ just enforces $G_i\lambda_x = 0$.) The barrier removes
the failure: substituting $s = \kappa/y_g$ (elementwise) gives
$W = y_g^2/(\kappa + \delta_\xi\gamma_\xi^{-1}y_g^2)$, smooth along the central
path, so $\nabla_\theta L$ is continuous across active-set changes~\cite{tracy}.

\subsection{Assembling the parameter gradients}
With $(\lambda_x,\lambda_f)$ from \eqref{eq:Wsys} (or \eqref{eq:J-adjoint}) and
$\bar\lambda_g$ from \eqref{eq:W}, the contraction
$\nabla_\theta L = (\partial F/\partial\theta)^\top\lambda$ of \eqref{eq:adjoint}
takes, at the NLP level, the compact form
\begin{equation}
\label{eq:nlp-grads}
\nabla_\theta L
= \bigl[\partial_\theta\nabla_x\mathcal L\bigr]^\top\lambda_x
+ \bigl[\partial_\theta f\bigr]^\top\lambda_f
+ \bigl[\partial_\theta (g - h)\bigr]^\top\bar\lambda_g ,
\end{equation}
three mixed vector--Jacobian products of the problem functions evaluated at
$(w^\star,\theta)$ --- the adjoint sensitivity of \cite[Eq.~(13)]{frey}. For
parameters that enter only the cost (the typical learned-MPC case),
$\partial_\theta f = 0$ and $\partial_\theta(g-h) = 0$, and \eqref{eq:nlp-grads}
collapses to the single contraction
$\nabla_\theta L = [\partial_\theta\nabla_x\ell(x^\star;\theta)]^\top\lambda_x$,
one vector--Jacobian product of the cost gradient.

\paragraph{QP-layer special case.} When the layer is the QP \eqref{eq:qp} itself,
with data $\theta = (P,q,C,c,G,h)$ and $H = P$, expanding \eqref{eq:nlp-grads}
datum-by-datum gives the familiar table
\begin{equation}
\label{eq:paramgrads}
\begin{aligned}
\nabla_P L &= \half\bigl(\lambda_x x^{\star\top} + x^\star\lambda_x^\top\bigr),
& \nabla_q L &= \lambda_x,\\
\nabla_C L &= \lambda_f x^{\star\top} + y_f\lambda_x^\top,
& \nabla_c L &= -\lambda_f,\\
\nabla_G L &= \bar\lambda_g x^{\star\top} + y_g\lambda_x^\top,
& \nabla_h L &= -\bar\lambda_g .
\end{aligned}
\end{equation}

\medskip
\noindent\textbf{Summary.} The forward pass drives each QP subproblem to the root
of \eqref{eq:pkkt} by a projection-free ADMM whose slack update is the closed-form
prox of the log barrier \eqref{eq:retr-update}, holding $s\circ y_g = \kappa\ones$
exactly at every iterate. The backward pass differentiates the \emph{NLP's}
smoothed KKT system \eqref{eq:nlp-pkkt} at the converged SQP solution with the
exact Lagrangian Hessian \eqref{eq:H}~\cite{frey}; the adjoint system
\eqref{eq:Wsys} shares the block-tridiagonal structure --- and the linear solver
--- of the forward pass. The parameters $(\kappa, \delta_\xi, \gamma_\xi)$
interpolate smoothly between this solver, TurboMPC's slack-smoothed ADMM, and
OSQP's projection ADMM, keeping the solution map differentiable for every
$\kappa > 0$.

\begin{thebibliography}{9}\footnotesize
\itemsep1pt

\bibitem{turbompc}
G.~Bravo-Palacios, J.~Zhang, Z.~Pestrikov, B.~Plancher, and T.~Lew.
\emph{TurboMPC: Fast, Scalable, and Differentiable Model Predictive Control on the
GPU}. arXiv:2606.24039, 2026.

\bibitem{osqp}
B.~Stellato, G.~Banjac, P.~Goulart, A.~Bemporad, and S.~Boyd.
\emph{OSQP: An operator splitting solver for quadratic programs}.
Mathematical Programming Computation, 12(4):637--672, 2020.

\bibitem{boyd}
S.~Boyd, N.~Parikh, E.~Chu, B.~Peleato, and J.~Eckstein.
\emph{Distributed optimization and statistical learning via the alternating
direction method of multipliers}.
Foundations and Trends in Machine Learning, 3(1):1--122, 2011.

\bibitem{parikh}
N.~Parikh and S.~Boyd.
\emph{Proximal algorithms}.
Foundations and Trends in Optimization, 1(3):127--239, 2014.

\bibitem{tracy}
K.~Tracy and Z.~Manchester.
\emph{On the differentiability of the primal-dual interior-point method}.
arXiv:2406.11749, 2024.

\bibitem{optnet}
B.~Amos and J.~Z. Kolter.
\emph{OptNet: Differentiable optimization as a layer in neural networks}.
ICML 2017, PMLR 70:136--145.

\bibitem{frey}
J.~Frey, K.~Baumg\"artner, G.~Frison, D.~Reinhardt, J.~Hoffmann, L.~Fichtner,
S.~Gros, and M.~Diehl.
\emph{Differentiable nonlinear model predictive control}.
arXiv:2505.01353, 2025.

\bibitem{nw}
J.~Nocedal and S.~J. Wright.
\emph{Numerical Optimization}, 2nd ed. Springer, 2006. (Chs.~16, 18, 19.)

\end{thebibliography}

\end{document}
